Calibration-set gains remain positive but decline as the calibration subsets grow.
At the smallest sizes, mean fitted gains range from 0.256 percentage points to 1.302 percentage points,
compared with 0.009 percentage points to 0.451 percentage points at the largest sizes. Test-set differences
generally move toward zero as calibration-set data increase, although the pattern
is not monotone for every dataset. At the largest sizes, mean test-set
accuracy differences are $-0.054$ percentage points on MMLU, $+0.052$ percentage points on
TriviaQA, $-0.078$ percentage points on MATH, $-0.095$ percentage points on SimpleQA, and $-0.025$ percentage points on
LiveCodeBench. The three-model class also spends more on average, with
budget-normalized cost differences of 0.21\% to 2.15\% at those sizes.
Common feasible coverage ranges from 99.32\% to 99.82\% across datasets.
These results suggest that more calibration-set data reduce losses from selecting
additional depth, but do not reveal a substantial average test-set advantage
from a third stage at the tested sizes. Persistent fitted gains, particularly
on MMLU and TriviaQA, do not transfer into comparable test-set gains.
